
% Resumo
\begin{resumo}[Resumo]

Apresente de forma concisa os pontos relevantes de seu trabalho, de forma que qualquer pessoa consiga ter uma visão rápida e clara de seu conteúdo e conclusões. O resumo é formado por apenas um parágrafo, sem tópicos, contendo, preferencialmente, entre 150 e 500 palavras. Convém evitar fórmulas e equações no resumo, pois, além de ser uma recomendação na NBR 6028 (Resumo), as expressões com caracteres especiais, sobretudo com subscritos ou sobrescritos, geralmente não são lidas adequadamente pelo site do Repositório Digital da UFPE – Attena, onde seu resumo e abstract ficarão disponíveis, o que dificulta também a recuperação pela BDTD Nacional, podendo prejudicar a visibilidade de seu trabalho. Nas palavras-chave indique de 3 a 6 termos que representem o conteúdo do trabalho, preferencialmente escolhidos do vocabulário controlado indicado pela biblioteca setorial de seu centro. Confira aqui qual vocabulário controlado é recomendado pela biblioteca de seu centro.


 \vspace{\onelineskip}
    
 \noindent
 \textbf{Palavras-chaves}:texto; texto; texto; texto; texto

 \textcolor{red}{RESUMO é um elemento OBRIGATORIO}

 
\end{resumo}


